\documentclass[11pt]{article}
\usepackage[T1]{fontenc}
\usepackage[margin=1in]{geometry}
\usepackage{amsmath,amssymb,amsthm,microtype,graphicx}
\usepackage[colorlinks=true,urlcolor=blue,citecolor=blue]{hyperref}
\newtheorem{theorem}{Theorem}
\newtheorem{corollary}[theorem]{Corollary}
\newtheorem{remark}[theorem]{Remark}
\DeclareMathOperator{\Ent}{Ent}
\DeclareMathOperator{\diam}{diam}
\title{Transport--entropy inequalities for Connes--Rieffel metrics}
\author{Solution packet for \texttt{fa\_banach\_001}}
\date{11 August 2026}

\begin{document}
\maketitle

\noindent\textbf{Status.}
Scoped full resolution of the transport question in the finite tracial
framework of Junge--Zeng: an exact necessary-and-sufficient exponential-moment
criterion for an arbitrary seminorm $L$, an automatic positive answer for
finite-diameter (hence compact) quantum metric spaces, and a sharp kernel
obstruction for general spectral triples.  This packet does not claim to
resolve the source's separate question about deriving $L_p$ Poincar\'e
inequalities from logarithmic Sobolev inequalities in full generality.

\section*{Source question}

Marius Junge and Qiang Zeng, \emph{Noncommutative martingale deviation and
Poincar\'e type inequalities with applications}, arXiv:1211.3209, published in
\emph{Probability Theory and Related Fields} 161 (2015), 449--507,
DOI \href{https://doi.org/10.1007/s00440-014-0552-1}{10.1007/s00440-014-0552-1}.
On page 4 of the cached source version, after defining
\[
 \rho_L(\phi,\psi)=\sup\{|\phi(a)-\psi(a)|:L(a)\leq1\}
\]
and noting that a spectral triple has $L(a)=\|[D,a]\|$, the authors write:
\begin{quote}
``It is an interesting question but beyond the scope of this paper to
determine whether the transportation inequality is possible for $\rho_L$.''
\end{quote}

\begin{center}
\includegraphics[width=0.96\linewidth]{figures/open_problem_crop.png}
\end{center}

\section*{Definitions and scope}

Let $(\mathcal M,\tau)$ be a finite von Neumann algebra with normal faithful
tracial state.  Let $\mathcal A_{\rm sa}\subseteq\mathcal M_{\rm sa}$ be a real
linear unital subspace, and let $L:\mathcal A_{\rm sa}\to[0,\infty)$ be a
seminorm with $L(1)=0$.  For a density $f\in L_1(\mathcal M,\tau)_+$ with
$\tau(f)=1$, write
\[
 \phi_f(a)=\tau(fa),\qquad
 \Ent_\tau(f)=\tau(f\log f),
\]
with the usual value $+\infty$ allowed.  Define
\[
 \rho_L(\phi_f,\tau)
 =\sup\bigl\{|\tau((f-1)a)|:
       a\in\mathcal A_{\rm sa},\ L(a)\leq1\bigr\}.
\]
The transportation inequality $T_1(C)$ means
\begin{equation}\label{eq:T1}
 \rho_L(\phi_f,\tau)\leq\sqrt{2C\,\Ent_\tau(f)}
 \quad\hbox{for every density }f.
\end{equation}
This is exactly the normal-state, finite-tracial setting used by the source.

\section*{Proof intuition}

The carré-du-champ does no work in the source's final Bobkov--G\"otze
duality argument: only the collection of admissible observables matters.
Replacing that collection by the $L$-unit ball therefore gives an exact
criterion for $\rho_L$.  Entropy is convex-dual to the logarithmic moment
generating function, so $T_1(C)$ is equivalent to a common subgaussian bound
for all centered observables with $L(a)\leq1$.

Finite metric diameter then supplies this bound without any semigroup or
curvature hypothesis.  Every admissible self-adjoint observable has spectral
range at most the diameter.  Tilting its spectral distribution and bounding
the tilted variance by one quarter of the squared range gives Hoeffding's
constant $\Delta_L^2/8$ in the logarithmic moment generating function, hence
$C=\Delta_L^2/4$.  In the opposite direction, a non-scalar element killed by
$L$ can be scaled at no cost and makes $\rho_L$ infinite against an explicit
bounded density.

\section*{Exact characterization and optimal constant}

Define
\begin{equation}\label{eq:CL}
 C_L=\sup_{\substack{a\in\mathcal A_{\rm sa},\ L(a)\leq1\\ t\in\mathbb R\setminus\{0\}}}
 \frac{2}{t^2}\log\tau\!\left(
  e^{t(a-\tau(a)1)}\right)\in[0,\infty].
\end{equation}

\begin{theorem}[Seminorm Bobkov--G\"otze duality]\label{thm:duality}
For $C\in[0,\infty)$, inequality \eqref{eq:T1} holds if and only if
\begin{equation}\label{eq:mgf}
 \log\tau\!\left(e^{t(a-\tau(a)1)}\right)\leq\frac{Ct^2}{2}
\end{equation}
for every $a\in\mathcal A_{\rm sa}$ with $L(a)\leq1$ and every
$t\in\mathbb R$.  Consequently $C_L$ is the optimal $T_1$ constant: a finite
constant exists exactly when $C_L<\infty$, and then \eqref{eq:T1} holds with
$C=C_L$.
\end{theorem}

\begin{proof}
Put $x=a-\tau(a)1$.  Since $L(1)=0$, the seminorm property gives
$L(x)=L(a)$.  The Gibbs variational formula in a finite tracial algebra is
\begin{equation}\label{eq:gibbs}
 \log\tau(e^h)
 =\sup_{\substack{f\geq0\\\tau(f)=1}}
   \{\tau(fh)-\Ent_\tau(f)\}
\end{equation}
for bounded self-adjoint $h$.  Indeed, the inequality ``$\geq$'' follows by
taking $f=e^h/\tau(e^h)$, while the reverse inequality is positivity of the
relative entropy of $f$ with respect to that Gibbs density.  Truncation gives
the same formulas for all finite-entropy $L_1$ densities needed below.

Assume \eqref{eq:mgf}.  Formula \eqref{eq:gibbs}, with $h=tx$, gives for
$t>0$
\[
 t\tau(fx)\leq \Ent_\tau(f)+\frac{Ct^2}{2}.
\]
Optimizing in $t$ yields
$\tau(fx)\leq\sqrt{2C\Ent_\tau(f)}$.  Applying this to $-x$ and taking the
supremum over the $L$-unit ball proves \eqref{eq:T1}.

Conversely, suppose \eqref{eq:T1}.  For $L(a)\leq1$ it implies
\[
 t\tau(fx)-\Ent_\tau(f)
 \leq t\sqrt{2C\Ent_\tau(f)}-\Ent_\tau(f)
 \leq \frac{Ct^2}{2}
\]
for $t\geq0$, because the supremum of
$t\sqrt{2Cs}-s$ over $s\geq0$ is $Ct^2/2$.  Taking the supremum over $f$ in
\eqref{eq:gibbs} proves \eqref{eq:mgf} for $t\geq0$; replacing $a$ by $-a$
handles $t<0$.  The statement about $C_L$ now follows directly from
\eqref{eq:CL}.
\end{proof}

\section*{Finite diameter gives a universal positive answer}

Let $S(\mathcal M)$ denote the state space and set
\[
 \Delta_L=\diam(S(\mathcal M),\rho_L)
 =\sup_{\phi,\psi\in S(\mathcal M)}\rho_L(\phi,\psi).
\]

\begin{theorem}[Diameter theorem]\label{thm:diameter}
If $\Delta_L<\infty$, then
\[
 C_L\leq\frac{\Delta_L^2}{4}
 \quad\hbox{and hence}\quad
 \rho_L(\phi_f,\tau)
 \leq \Delta_L\sqrt{\frac{\Ent_\tau(f)}{2}}
\]
for every density $f$.
\end{theorem}

\begin{proof}
Fix $a=a^*$ with $L(a)\leq1$.  If $m=\min\sigma(a)$ and
$M=\max\sigma(a)$, evaluation on states attaining the spectral extrema gives
\[
 M-m=\sup_{\phi,\psi\in S(\mathcal M)}|\phi(a)-\psi(a)|\leq\Delta_L.
\]
Let $x=a-\tau(a)1$ and $H(t)=\log\tau(e^{tx})$.  Functional calculus reduces
this to the logarithmic moment generating function of the scalar spectral
measure of $x$ under $\tau$.  Thus $H(0)=H'(0)=0$, and
\[
 H''(t)=\operatorname{Var}_{\mu_t}(x)leq\frac{(M-m)^2}{4}
 \leq\frac{\Delta_L^2}{4},
\]
where $\mu_t$ is the exponentially tilted spectral measure.  The variance
bound follows from
$\operatorname{Var}(X)\leq(M-\mathbb EX)(\mathbb EX-m)\leq(M-m)^2/4$
for every random variable supported in $[m,M]$.  Integrating twice, for
positive or negative $t$, gives
\[
 \log\tau(e^{tx})=H(t)\leq\frac{\Delta_L^2t^2}{8}.
\]
Theorem \ref{thm:duality} applies with $C=\Delta_L^2/4$.
\end{proof}

\begin{corollary}[Compact quantum metric spaces]\label{cor:compact}
Suppose $L$ is a Rieffel Lip-norm on the self-adjoint part of a unital
$C^*$-algebra represented in $(\mathcal M,\tau)$, so that $\rho_L$ metrizes
the weak-star topology on the state space.  Then $\Delta_L<\infty$ and the
transportation inequality in Theorem \ref{thm:diameter} holds.  In particular,
the answer is positive for every finite-tracial spectral triple for which
$L(a)=\|[D,a]\|$ is a Lip-norm.
\end{corollary}

\begin{proof}
The state space is weak-star compact.  If $\rho_L$ metrizes that topology, it
is a compact metric space and therefore has finite diameter.  Apply Theorem
\ref{thm:diameter}.
\end{proof}

\section*{The kernel obstruction}

\begin{theorem}[Failure for degenerate spectral triples]\label{thm:kernel}
If there is a non-scalar $a=a^*\in\mathcal A_{\rm sa}$ with $L(a)=0$, then
there exists a bounded, strictly positive density $f$ with finite entropy such
that
\[
 \rho_L(\phi_f,\tau)=\infty.
\]
Consequently no finite-constant transportation inequality can hold.
\end{theorem}

\begin{proof}
Set $b=a-\tau(a)1$.  Then $b\neq0$, $\tau(b)=0$, and $L(b)=0$.  Choose a
nonzero real $\varepsilon$ with
$|\varepsilon|<\|b\|_\infty^{-1}$ and put $f=1+\varepsilon b$.  Then $f$ is a
bounded strictly positive density, so $\Ent_\tau(f)<\infty$.  For every
$s\in\mathbb R$, the observable $sb$ is admissible because $L(sb)=0$, while
\[
 |\tau((f-1)sb)|=|s\varepsilon|\,\tau(b^2).
\]
This tends to infinity with $|s|$, proving the claim.
\end{proof}

\section*{Verification and relation to the source}

The proof has four independent checks.
\begin{enumerate}
\item The two directions of Theorem \ref{thm:duality} use the same Gibbs
variational identity, so the constants can be checked by the scalar maximization
$\sup_{s\geq0}(t\sqrt{2Cs}-s)=Ct^2/2$.
\item The diameter constant is checked by the extremal two-point distribution:
the universal centered log-moment bound for range length $\Delta$ is
$\Delta^2t^2/8$, exactly corresponding to $C=\Delta^2/4$.
\item The kernel witness is explicit, positive, normalized, and has finite
entropy; scaling the zero-seminorm observable forces infinite distance.
\item For a Rieffel Lip-norm, compactness of the state space immediately gives
finite diameter, so no unstated semigroup, curvature, or diffusion hypothesis
enters Corollary \ref{cor:compact}.
\end{enumerate}

Junge--Zeng prove a Bobkov--G\"otze equivalence for their specific
carr\'e-du-champ unit ball (their Theorem 3.17) and a finite-diameter criterion
for that same semigroup metric (their Corollary 3.19).  The new step here is to
observe and prove that the variational mechanism is completely independent of
the carr\'e du champ, apply it to the separately posed Connes--Rieffel metric,
identify the optimal constant \eqref{eq:CL}, obtain the diameter constant
$\Delta_L^2/4$, and pair it with the exact kernel obstruction.

\section*{Novelty search and limitations}

The run's solution, attempt, registry, and literature indexes were searched for
arXiv:1211.3209 and the phrases ``Rieffel metric'', ``Connes metric'',
``transport entropy'', and ``spectral triple''; no duplicate result was found.
A bounded arXiv search used the combinations ``Rieffel metric transport
entropy inequality spectral triple'', ``transportation inequality compact
quantum metric'', and ``rho\_L entropy transportation inequality''.  It found
work establishing that various Connes seminorms are compact quantum metrics,
including Rieffel's group $C^*$-algebra examples and later spectral-metric
examples, but no exact later paper resolving the source's entropy inequality.

The result is mathematically elementary once the correct dual formulation is
isolated, so originality confidence is lower than validity confidence; a
reviewer should run a broader citation search before public novelty claims.
The theorem concerns normal states with densities relative to a finite trace,
exactly the source's probabilistic framework.  It does not define entropy for
arbitrary singular states, does not assert finite diameter for every spectral
triple, and does not resolve the distinct LSI-to-$L_p$ Poincar\'e question.

\section*{Human-review recommendation}

\textbf{Likely valid; promote after checking scope and novelty.}  The formal
proof is short and constant-exact.  The most important review points are
(i) retaining the finite-tracial normal-state scope in any public statement,
(ii) confirming that the intended meaning of ``possible'' in the source
accepts an exact criterion plus the compact-quantum-metric corollary, and
(iii) checking whether the seminorm-agnostic abstraction of the source's
Theorem 3.17 has appeared elsewhere.

\begin{thebibliography}{9}
\bibitem{JZ}
M.~Junge and Q.~Zeng,
\emph{Noncommutative martingale deviation and Poincar\'e type inequalities
with applications}, Probab. Theory Related Fields 161 (2015), 449--507,
arXiv:1211.3209.

\bibitem{Rie}
M.~A. Rieffel,
\emph{Compact quantum metric spaces}, in Operator algebras, quantization, and
noncommutative geometry, Contemp. Math. 365 (2004), 315--330.

\bibitem{RieGroup}
M.~A. Rieffel,
\emph{Group $C^*$-algebras as compact quantum metric spaces},
Doc. Math. 7 (2002), 605--651, arXiv:math/0205195.

\bibitem{AC}
C.~Antonescu and E.~Christensen,
\emph{Group $C^*$-algebras, metrics and an operator theoretic inequality},
J. Funct. Anal. 214 (2004), 247--259, arXiv:math/0211312.
\end{thebibliography}

\end{document}
